\section{Discussion}
\label{sec:discussion}

\subsection*{What the results mean, at the level the evidence supports}
\TODO{First paragraph restates the headline finding in one sentence and immediately
bounds it. Resist writing "cheetah movement" anywhere in this section.}

\subsection{Where permeability has declined, and what co-occurred}
\TODO{H1. Frame as spatial co-occurrence, not causation (CAU-014). If the dominant
pressure differs by country, report the heterogeneity rather than compressing it into
a single anthropogenic narrative.}

\subsection{The protection gap as a governance question}
\TODO{H2. If the outside-PA share is high, the defensible reading is that
connectivity for this species depends on land under other management, which is a
statement about where conservation effort must operate, not evidence that protected
areas have failed. Note the confound that protected areas are often sited in land of
low agricultural value.}

\subsection{Robust versus model-sensitive linkages}
\TODO{H3. This is the paper's most transferable contribution. Robust locations
support action; model-sensitive locations define the survey agenda. Note the caveat
that scenarios sharing input layers can agree artificially.}

\subsection{From uncertainty to a survey design}
\TODO{H4. The monitoring index is decision support under stated assumptions, not an
objective ranking (CAU-020). Say what a partner organisation would actually do with
it, and acknowledge that accessibility and cost will reorder the list in practice.}

\subsection{Comparison with regional connectivity work}
\TODO{Position against the elephant and multispecies transfrontier assessments
\parencite{naidoo2024kaza,brennan2020multispecies} as methodological neighbours, and
against the fencing-effects literature \parencite{osipova2018fencing}. Be explicit
that those are different species with different movement ecology: the methods
transfer, the resistance values do not (CAU-015).}
